\scanpage{2}
\mspage{1}
Réalisations topologiques des réseaux \hfill (Juin 86)

\margin{\unsure{CF II|Ch II}}

1. Ficelle topologique. (6 Juin)

Soit $X$ espace top., $A\subset X$ partie fermée \inserted{non vide} de $X$.
\deleted{\unsure{Soit}} $X/A$ l'espace déduit de $X$ en contractant $A$ en un pt,
\inserted{$a$ le $p^t$ déduit de $A$ par contraction.}
Si $X'$ est une partie de $X$ contenant $A$, alors $X'/A\hookrightarrow X/A$ identifie $X'/A$ à un sous-espace top. de $X$
[les fermés de $X'/A$ s'identifient aux fermés \inserted{$Y'$} \deleted{\unsure{$Y$}} de $X'$ qui ou bien contiennent $A$, ou bien ne le rencontrent pas (i.e. tels que $Y'\cap A\in\{\varnothing,A\}$), et ce sont les traces sur $X'$ des fermés $Y$ de $X$ qui ont ces \unsure{$=$s|$\Longrightarrow$} propriétés, lesquels correspondent \inserted{d'ailleurs} aux fermés de $X/A$ \ldots]

On dit que $A$ est \uline{régulièrement immergé} dans $X$
si $\exists$ voisinage \deleted{\unsure{$V$}} conique de $a$ dans $X/A$
i.e. voisinage \inserted{\uline{fermé}}
$V\overset{h}{\simeq}\operatorname{\text{Cône}}(\unsure{B|S})$
$\unsure{B|S}$ espace topologique compact, $h(a)=o$.
(centre du cône),
\inserted{et $V^\circ=V\setminus B$ ouvert dans $X/A$.}
\deleted{Sur $V$} Alors \deleted{\unsure{le}} sur $V$ il y a (via $h$) une structure conique et en part. considérons les
\[
tV=V_t\qquad 0\le t\le1
\]
qui sont, pour $t\in]0,1]$, homéom. à $V$, et forment

\margin{\unsure{Souvent dit} ``\unsure{saturé}'' \deleted{\unsure{dans}} \unsure{d'après}
\deleted{\unsure{NB} \unsure{d'usage} $\unsure{\partial V|V^\circ}$ \illegible{2}}}

\margin{\uline{NB} $\dot V=B$\quad $\dot V_t=B_t$}
\mspage{1}

\rotatebox{90}{\unsure{EXTRA STRONG}}
\rotatebox{90}{\unsure{EXTRA STRONG}}
